\chapter{Lexical Structure}
\section{Character Set}
EDADSL uses ASCII for identifiers, keywords, operators, and grammar symbols. Unicode characters are permitted within string literals and comments.
\section{Whitespace}
Whitespace is defined using the following Regular Expression:
\begin{verbatim}
RegEx:
Whitespace                    = [ \t\n\r]+
\end{verbatim}

\section{Comments}
Comments are defined using the following Regular Expressions:
\begin{verbatim}
RegEx:
Line_Comments                 = "#c#"[^\n]*
Multiline_Comments            = "#c#", "{", [any characters], "}#c#"
\end{verbatim}

Multiline comments start with \texttt{\#c\#\{} and end with \texttt{\}\#c\#}. Nested multiline comments are not supported.

Examples:
\begin{verbatim}
#c# This is a line comment
$x [int]= 5  #c# Inline comment

#c#{
This is a
multiline comment
}#c#

$x [int]= 5 #c#{ inline }#c# + 3
\end{verbatim}

\section{Identifiers}
Identifiers are defined using the following Regular Expressions:
\begin{verbatim}
RegEx:
id                            = [_a-zA-Z][_a-zA-Z0-9]*
\end{verbatim}

\section{Keywords}
The language uses the following Keywords defined using the following Regular Expressions:
\begin{verbatim}
Generic_Keywords:    start halt if else otherwise for while do break continue
                     gotomarker gotojump 
                     writetolog make.pcb make.schematic export 
                     echo return let letop skip function L::
                     null assert

Physical_Keywords:   phys\.[a-z]+
                     Known Physical Keywords:
                     phys.prox  phys.creep  phys.layer
                     phys.shrinknet  phys.mkplane
                     phys.viafill  phys.connectioncurrent
                     phys.align  phys.connmatch  phys.board
                     phys.mktext

Electrical_Keywords: elec\.[a-z]+
                     Known Electrical Keywords:
                     elec.part  elec.connect  elec.net
                     elec.tag  elec.untag

System_Keywords:     syst\.[a-z]+
                     Known System Keywords:
                     syst.flag\.[a-z]+
\end{verbatim}
Keywords are reserved identifiers and may not be used as variable names, function names, or constant names. Keywords are case-sensitive: \texttt{if} is a keyword, but \texttt{If} or \texttt{IF} is an identifier.

\section{Literals}
For the Lexical Structure of Literals see:  Section~\ref{sec:generic-types} (Types).

\section{Operators}
For the Lexical Structure of built in Operators see:  Section~\ref{sec:operators} (Operators). \\
User-defined Operators are defined using the following Regular Expressions:
\begin{verbatim}
User_Operators     = o\.u\.[_a-zA-Z0-9]+
\end{verbatim}
The \texttt{o.u.} prefix is distinct from dot-access chains (\texttt{obj.property}): \texttt{o.u.} is always followed by an operator name and used in infix position, while dot-access is always \texttt{identifier.identifier} and denotes property access.

\section{Delimiters}
EDADSL uses the following Delimiters:
\begin{verbatim}
() [] {} . , : ; @ ~ | %N
\end{verbatim}

The \texttt{\%N} token is a statement separator, not a comment. It allows multiple statements on a single line. The \texttt{::} token denotes a code block opener. The semicolon \texttt{;} appears in modifier syntax (\texttt{;M=}, \texttt{;P=}, \texttt{;T=}) and property separators.

\section{Tokenization Rules}
The Tokenization follows the following Rules: \\
The Lexer uses the Longest Match Rule: in Cases where multiple valid tokenizations are possible the longest token is chosen. \\
If Lexeme matches both a Keyword and an Identifier it is classified as a Keyword. \\
Whitespace and Comments are discarded.

\section{Statement Separators}
Statements within a code block are separated by implicit newlines. The \texttt{\%N} token acts as an explicit statement separator, allowing multiple statements on a single line:
\begin{verbatim}
#c# These are equivalent:
::{
    f.print("a")
    f.print("b")
}

::{
    f.print("a") %N f.print("b")
}
\end{verbatim}
The \texttt{\%N} token is primarily useful inside inline code blocks where writing multiple lines is impractical.

\section{End of File}
Source files must end with a newline character. A line comment at the end of the file without a trailing newline is permitted but discouraged.

